\documentclass[12pt,a4paper]{ctexart}

% === 字体设置 ===
\setCJKmainfont{SimSun}
\setmainfont{Times New Roman}

% === 宏包 ===
\usepackage{amsmath}
\usepackage{amssymb}
\usepackage[version=4]{mhchem}
\usepackage{graphicx}
\usepackage{paracol}
\usepackage{xcolor}
\usepackage{fancyhdr}
\usepackage{geometry}

% === 页面设置 ===
\geometry{
  a4paper,
  left=1.5cm,
  right=1.5cm,
  top=2cm,
  bottom=2cm,
  columnsep=1cm
}

% === 内部辅助：计算 paracol 列宽 ===
\newlength{\colwidthinner}
\newcommand{\calccolwidth}{%
  \setlength{\colwidthinner}{\dimexpr\linewidth-2\fboxsep-2\fboxrule\relax}%
}

% 右栏：修改建议框
\newcommand{\correctionbox}[1]{%
  \calccolwidth
  \colorbox{blue!8}{\parbox{\colwidthinner}{#1}}%
}

% === 页眉页脚 ===
\pagestyle{fancy}
\fancyhf{}
\fancyhead[L]{校对报告}
\fancyhead[R]{\leftmark}
\fancyfoot[C]{\thepage}

\begin{document}

\title{test}
\author{}
\date{}
\maketitle

\begin{paracol}{2}

**一本班1**（多选）\\
如图所示\textsuperscript{\textcolor{red}{\textcircled{1}}}，利用斜面从货车上卸货，每包货物的质量$m=20\mathrm{kg}$，斜面倾角$\alpha ={37}^{\circ }$，斜面的长度$l=0.5\mathrm{m}$，货物与斜面间的动摩擦因数$\mu =0.2$，货物从斜面顶端滑到底端的过程中，有（ ）（取$g=10\mathrm{m}/{\mathrm{s}}^{2}$）\textsuperscript{\textcolor{red}{\textcircled{2}}}\\
A.                                  重力对货物做的功为$60\mathrm{J}$\\
B.                                  斜面支持力对货物做的功为$0$\\
C.                                  斜面摩擦力对货物做的功为$16\mathrm{J}$\\
D.                                  合外力对货物做的功为$76\mathrm{J}$\\
答案:\\
AB\\
解答：                             重力做功：${W}_{\mathrm G}=mgL\sin {37}^{\circ }=60\mathrm{J}$\textsuperscript{\textcolor{red}{\textcircled{3}}}，选项$\mathrm{A}$正确；\\
支持力方向与位移方向垂直，不做功，则有：${W}_{\mathrm N}=0$，选项$\mathrm{B}$正确；\\
滑动摩擦力为：$f=\mu {F}_{\mathrm{N}}=0.2\times 20\times 10\times 0.8\mathrm{N}=32\mathrm{N}$，则摩擦力做功为：${W}_{\mathrm{f}}=-fL=-32\times 0.5\mathrm{J}=-16\mathrm{J}$\textsuperscript{\textcolor{red}{\textcircled{4}}}，选项$\mathrm{C}$错误；\\
合外力做功为：${W}_{合}={W}_{\mathrm{N}}+{W}_{\mathrm{f}}+{W}_{\mathrm{G}}=60\mathrm{J}-16\mathrm{J}=44\mathrm{J}$，选项$\mathrm{D}$错误。\\

\switchcolumn

\correctionbox{\textcircled{1} 题干提及“如图所示”但未提供对应配图，存在图片缺失，需补充配图以完善题干表述。}
\medskip

\correctionbox{\textcircled{2} 调整重力加速度说明的位置至问题前，补充多选题常规引导语，使题干表述通顺、逻辑清晰。}
\medskip

\correctionbox{\textcircled{3} 与题干中斜面长度的符号\$l\$保持一致，规范物理量符号使用。}
\medskip

\correctionbox{\textcircled{4} 与题干中斜面长度的符号\$l\$保持一致，规范物理量符号使用。

---}
\medskip

\switchcolumn*

\end{paracol}

\end{document}
